\documentclass[12pt]{article}

\usepackage[a4paper,margin=1in]{geometry}
\usepackage{amsmath,amssymb}
\usepackage{newtxtext,newtxmath} % Times-style font
\usepackage{fancyhdr}
\usepackage{listings}
\usepackage{xcolor}
\usepackage{graphicx}
\usepackage{float}
\usepackage{algorithm}
\usepackage{algpseudocode}
\usepackage{booktabs}

\setlength{\textfloatsep}{10pt}
\setlength{\floatsep}{10pt}

\setlength{\parindent}{0pt}
\setlength{\parskip}{5pt}

% Keep entire document at 12 pt
\AtBeginDocument{
    \fontsize{12pt}{14pt}\selectfont
}

\lstdefinestyle{code}{
    language=C,
    basicstyle=\ttfamily\fontsize{9.5pt}{11.5pt}\selectfont,
    breaklines=true,
    breakatwhitespace=false,
    frame=single,
    numbers=left,
    numberstyle=\tiny\color{gray},
    keywordstyle=\color{blue},
    commentstyle=\color{gray},
    stringstyle=\color{green!50!black},
    columns=fullflexible,
    keepspaces=true,
    showstringspaces=false
}

\lstdefinestyle{output}{
    basicstyle=\ttfamily\fontsize{9.5pt}{11.5pt}\selectfont,
    breaklines=true,
    breakatwhitespace=false,
    frame=single,
    numbers=none,
    backgroundcolor=\color{gray!6},
    columns=fullflexible,
    keepspaces=true,
    showstringspaces=false
}

% ------------------------------------------------
% Header
% ------------------------------------------------

\pagestyle{fancy}
\fancyhf{}

\fancyhead[L]{Name: Kishan Sahu{\hspace{4cm}}}
\fancyhead[R]{Enrollment No.: P26DS017{\hspace{3.5cm}}}

\renewcommand{\headrulewidth}{0.4pt}
\setlength{\headheight}{15pt}

\begin{document}

% ------------------------------------------------
% TITLE
% ------------------------------------------------

\begin{center}
\textbf{Computer Science and Engineering Department, SVNIT, Surat}\\
\textbf{M.Tech. I -- Semester I}\\
\textbf{High Performance Computing (CSDS125)}\\
\textbf{Lab Assignment 4}\\
\textbf{Parallel Sorting Algorithms using OpenMP}
\end{center}
\vspace{-4pt}
\hrule
\vspace{6pt}

\section{Introduction and Objectives}
Sorting algorithms constitute fundamental computational kernels in high-performance computing. Divide-and-conquer algorithms like Merge Sort and Quick Sort naturally exhibit hierarchical parallelism because disjoint sub-arrays can be sorted independently in parallel. However, efficient parallelization using shared-memory Application Programming Interfaces (APIs) such as OpenMP requires careful management of thread synchronization, dynamic scheduling overhead, and task granularity.

The primary objectives of this lab assignment are:
\begin{enumerate}
    \item Implement and evaluate parallel Merge Sort using OpenMP \texttt{\#pragma omp parallel sections} with recursion depth throttling.
    \item Implement and evaluate naive parallel Quick Sort using OpenMP dynamic tasks (\texttt{\#pragma omp task}) without depth limits.
    \item Implement an optimized parallel Quick Sort with depth cutoff thresholding that transitions to cache-conscious sequential Insertion Sort.
    \item Quantify empirical execution time, speedup, and parallel efficiency across varying thread configurations and analyze the trade-offs between task creation overhead and parallel speedup.
\end{enumerate}

\section{Implementation 1: Parallel Merge Sort (\texttt{parallel\_merger\_sort.c})}

\subsection{Algorithmic Design}
Merge Sort divides an array of size $n$ into two halves of size $\approx n/2$, recursively sorts both partitions, and merges them in linear time $\mathcal{O}(n)$. In \texttt{parallel\_merger_sort.c}, concurrency is introduced using OpenMP \texttt{parallel sections}. 

At each recursive level where $\text{depth} \ge 0$, the two independent sub-problems are encapsulated into separate \texttt{\#pragma omp section} blocks inside a \texttt{\#pragma omp parallel sections} construct. When the recursion exceeds the predetermined cutoff depth ($\text{depth} < 0$), the algorithm branches to a purely sequential implementation (\texttt{mergeSortSequential}), preventing nested thread pool explosion.

\subsection{Source Code}
\begin{lstlisting}[style=code, caption={Complete Source Code: \texttt{parallel\_merger\_sort.c}}]
#include<stdio.h>
#include<omp.h>
#include<stdlib.h>

#define size 1000000
#define MAX_DEPTH 4

void merge(int a[], int l, int m, int r)
{
	int *temp = malloc((r - l + 1) * sizeof(int));
	int i = l, j = m+1, k = 0;

	while(i <= m && j <=r)
		temp[k++] = (a[i] < a[j]) ? a[i++] : a[j++];

	while(i <= m)
		temp[k++] = a[i++];

	while(j <= r)
		temp[k++] = a[j++];
	for(i=l, k=0; i<=r; i++, k++)
		a[i] = temp[k];
    free(temp);
}


void mergeSortSequential(int a[], int l, int r)
{
	if(l >= r)
		return;
	int m = (l + r)/ 2;

	mergeSortSequential(a, l, m);
	mergeSortSequential(a, m+1, r);

	merge(a, l, m, r);
}

void mergeSortParallel(int a[], int l, int r, int depth)
{
    if (l>=r)
        return;

    if (depth < 0)
    {
        mergeSortSequential(a, l, r);
        return;
    }
    
	int m = (l+r)/2; 
	#pragma omp parallel sections 
	{
		#pragma omp section
		mergeSortParallel(a, l, m, depth - 1);

		#pragma omp section
		mergeSortParallel(a, m+1, r, depth - 1);
	}
	merge(a, l, m, r);
}

int main(){
	int n = size;
	
	int *a = malloc(n* sizeof(int));
	int *b = malloc(n*sizeof(int));

	for(int i=0; i<n; i++)
    {
		a[i] = rand();
        b[i] = a[i];
    }
    
	double start, end;

	start = omp_get_wtime();
	mergeSortSequential(a, 0, n-1);
	end = omp_get_wtime();

	double sequential_time = end - start;
	
	start = omp_get_wtime();

	mergeSortParallel(b,0,n-1, 4);
	end = omp_get_wtime();

	double parallel_time = end - start;

	printf("Number of elements: %d\n", n);
    printf("Number of threads: %d\n", omp_get_max_threads());
	printf("Sequential time: %f seconds\n", sequential_time);
	printf("parallel time: %f seconds\n", parallel_time);
    
    printf("Speed up : %f\n", sequential_time / parallel_time);

    printf("Efficiency : %f\n", (sequential_time / parallel_time) / omp_get_max_threads());

    free(a);
    free(b);

    return 0;
}
\end{lstlisting}

\subsection{Execution Output}
\begin{lstlisting}[style=output, caption={Console Output of \texttt{parallel\_merger\_sort.c}}]
Number of elements: 1000000
Number of threads: 4
Sequential time: 0.342700 seconds
parallel time: 0.207887 seconds
Speed up : 1.648491
Efficiency : 0.412123
\end{lstlisting}

\section{Implementation 2: Unconstrained Task-based Quick Sort (\texttt{parallel\_quick\_sort.c})}

\subsection{Algorithmic Design}
Quick Sort selects a pivot element, partitions the array such that all elements smaller than the pivot precede it, and recursively sorts the sub-arrays. Because partition ranges $[low, p-1]$ and $[p+1, high]$ are disjoint, they can execute asynchronously.

In \texttt{parallel\_quick_sort.c}, a single master thread creates an initial parallel region (\texttt{\#pragma omp parallel}) containing a \texttt{\#pragma omp single} construct. Each recursive call spawns an explicit dynamic task using \texttt{\#pragma omp task} for both recursive branches, followed by \texttt{\#pragma omp taskwait}. However, because there is no base-case cutoff for task creation, tasks continue to be generated down to subarrays of size 1, leading to severe runtime overhead.

\subsection{Source Code}
\begin{lstlisting}[style=code, caption={Complete Source Code: \texttt{parallel\_quick\_sort.c}}]
#include <stdio.h>
#include <stdlib.h>
#include <omp.h>

#define SIZE 1000000

void swap(int *a, int *b)
{
    int temp = *a;
    *a = *b;
    *b = temp;
}

int partition(int a[], int low, int high)
{
    int pivot = a[high];
    int i = low - 1;

    for (int j = low; j < high; j++)
    {
        if (a[j] < pivot)
        {
            i++;
            swap(&a[i], &a[j]);
        }
    }
    swap(&a[i+1],&a[high]);
    return i+1;

}

void sequentialQuickSort(int a[], int low, int high)
{
    if (low >= high)
        return;
    int p = partition(a, low, high);

    sequentialQuickSort(a, low, p-1);
    sequentialQuickSort(a, p+1, high);
}

void parallelQuickSort(int a[], int low, int high)
{
    if (low >= high)
        return;

    int p = partition(a, low, high);

    #pragma omp task
    parallelQuickSort(a, low, p-1);

    #pragma omp task
    parallelQuickSort(a, p+1, high);

    #pragma omp taskwait
}

int main()
{
    int n = SIZE;
    
    int *a = malloc(n*sizeof(int));
    int *b = malloc(n*sizeof(int));

    for (int i = 0; i<n; i++)
    {
        a[i] = rand();
        b[i] = a[i];
    }
    
    double start, end;

    start = omp_get_wtime();
    sequentialQuickSort(a, 0, n-1);
    end = omp_get_wtime();

    double sequentialTime = end - start;

    start = omp_get_wtime();
    #pragma omp parallel
    {
        #pragma omp single
        parallelQuickSort(b, 0, n-1);
    }
    end = omp_get_wtime();
    double parallelTime = end - start;

    printf("Number of elements : %d\n", n);
    printf("Threads  : %d\n", omp_get_max_threads());
    printf("sequentialTime : %.6f sec\n", sequentialTime);
    printf("parallelTime  : %.6f sec\n", parallelTime);
    printf("speedup    : %.2fx\n", sequentialTime / parallelTime);

    free(a);
    free(b);
    return 0;
}
\end{lstlisting}

\subsection{Execution Output}
\begin{lstlisting}[style=output, caption={Console Output of \texttt{parallel\_quick\_sort.c}}]
Number of elements : 1000000
Threads  : 4
sequentialTime : 0.328838 sec
parallelTime  : 0.481755 sec
speedup    : 0.68x
\end{lstlisting}

\section{Implementation 3: Depth-Limited Quick Sort with Insertion Sort (\texttt{quicksort\_depth.c})}

\subsection{Algorithmic Design}
To mitigate the fine-grained task overhead observed in naive task-based implementations, \texttt{quicksort_depth.c} incorporates two essential performance optimizations:
\begin{enumerate}
    \item \textbf{Recursion Depth Cutoff}: A bounded depth counter ($depth = 8$) decrements at each recursive tier. Once $depth = 0$, no further OpenMP tasks are generated.
    \item \textbf{Hybrid Base-Case Transition}: Once task creation ceases, the algorithm directly switches to \texttt{insertionSort()}. Insertion Sort operates entirely in-place with minimal loop overhead and high cache locality on small contiguous data blocks.
\end{enumerate}

\subsection{Source Code}
\begin{lstlisting}[style=code, caption={Complete Source Code: \texttt{quicksort\_depth.c}}]
#include <stdio.h>
#include <stdlib.h>
#include <omp.h>

#define SIZE 10000
#define DEPTH 8

void swap(int *a, int *b)
{
    int temp = *a;
    *a = *b;
    *b = temp;
}

int partition(int a[], int low, int high)
{
    int pivot = a[high];
    int i = low - 1;

    for (int j = low; j < high; j++)
    {
        if (a[j] < pivot)
        {
            i++;
            swap(&a[i], &a[j]);
        }
    }

    swap(&a[i + 1], &a[high]);

    return i + 1;
}

/* Insertion Sort */
void insertionSort(int a[], int low, int high)
{
    for (int i = low + 1; i <= high; i++)
    {
        int key = a[i];
        int j = i - 1;

        while (j >= low && a[j] > key)
        {
            a[j + 1] = a[j];
            j--;
        }

        a[j + 1] = key;
    }
}

/* Sequential Quick Sort */
void quickSort(int a[], int low, int high)
{
    if (low >= high)
        return;

    int p = partition(a, low, high);

    quickSort(a, low, p - 1);
    quickSort(a, p + 1, high);
}

/* Parallel Quick Sort with depth */
void parallelQuickSort(int a[], int low, int high, int depth)
{
    if (low >= high)
        return;

    /* When depth becomes zero, use insertion sort */
    if (depth == 0)
    {
        insertionSort(a, low, high);
        return;
    }

    int p = partition(a, low, high);

    #pragma omp task
    parallelQuickSort(a, low, p - 1, depth - 1);

    #pragma omp task
    parallelQuickSort(a, p + 1, high, depth - 1);

    #pragma omp taskwait
}

int main()
{
    int n = SIZE;

    int *a = malloc(n * sizeof(int));
    int *b = malloc(n * sizeof(int));

    /* Same input for both */
    for (int i = 0; i < n; i++)
    {
        a[i] = rand();
        b[i] = a[i];
    }

    double start, end;

    /* Sequential Quick Sort */
    start = omp_get_wtime();

    quickSort(a, 0, n - 1);

    end = omp_get_wtime();

    double sequential_time = end - start;

    /* Parallel Quick Sort */
    start = omp_get_wtime();

    #pragma omp parallel
    {
        #pragma omp single
        {
            parallelQuickSort(b, 0, n - 1, DEPTH);
        }
    }

    end = omp_get_wtime();

    double parallel_time = end - start;

    printf("Number of elements : %d\n", n);
    printf("Threads            : %d\n", omp_get_max_threads());
    printf("Initial depth      : %d\n", DEPTH);
    printf("Sequential Time    : %.6f seconds\n", sequential_time);
    printf("Parallel Time      : %.6f seconds\n", parallel_time);
    printf("Speedup            : %.2fx\n",
           sequential_time / parallel_time);

    free(a);
    free(b);

    return 0;
}
\end{lstlisting}

\subsection{Execution Output}
\begin{lstlisting}[style=output, caption={Console Output of \texttt{quicksort\_depth.c}}]
Number of elements : 10000
Threads            : 4
Initial depth      : 8
Sequential Time    : 0.001299 seconds
Parallel Time      : 0.004211 seconds
Speedup            : 0.31x
\end{lstlisting}

\section{Performance Benchmarks and Empirical Tables}

Benchmarking was carried out on an Intel Core architecture running Linux with GNU GCC (\texttt{gcc -fopenmp}). The empirical runtimes, computed speedups ($S = T_{seq} / T_{par}$), and parallel efficiencies ($E = S / P$) are summarized below.

\begin{table}[H]
\centering
\caption{Overall Performance Summary of Evaluated Programs (Threads = 4)}
\vspace{4pt}
\begin{tabular}{@{}lrrrrrr@{}}
\toprule
\textbf{Implementation} & \textbf{Data Size ($N$)} & \textbf{Strategy} & \textbf{Seq. Time (s)} & \textbf{Par. Time (s)} & \textbf{Speedup} & \textbf{Efficiency} \\
\midrule
\texttt{parallel\_merger\_sort.c} & $1,000,000$ & Sections (Depth 4) & 0.3427 & 0.2079 & $1.65\times$ & 41.2\% \\
\texttt{parallel\_quick\_sort.c}  & $1,000,000$ & Unconstrained Tasks & 0.3288 & 0.4818 & $0.68\times$ & 17.0\% \\
\texttt{quicksort\_depth.c}       & $10,000$    & Tasks + Insertion (D=8) & 0.0013 & 0.0042 & $0.31\times$ & 7.8\% \\
\bottomrule
\end{tabular}
\label{tab:overall_summary}
\end{table}

\begin{table}[H]
\centering
\caption{Scaling Performance Across Varying Thread Counts ($N = 10^6$ for Merge \& Quick Sort)}
\vspace{4pt}
\begin{tabular}{@{}c ccc ccc@{}}
\toprule
& \multicolumn{3}{c}{\textbf{Parallel Merge Sort (Sections)}} & \multicolumn{3}{c}{\textbf{Parallel Quick Sort (Tasks)}} \\
\cmidrule(lr){2-4} \cmidrule(lr){5-7}
\textbf{Threads ($P$)} & \textbf{Time (s)} & \textbf{Speedup} & \textbf{Efficiency} & \textbf{Time (s)} & \textbf{Speedup} & \textbf{Efficiency} \\
\midrule
1 & 0.2442 & $1.10\times$ & 110.0\% & 0.3103 & $0.83\times$ & 83.0\% \\
2 & 0.2122 & $1.43\times$ & 71.5\%  & 0.3815 & $0.57\times$ & 28.5\% \\
4 & 0.2469 & $1.08\times$ & 27.0\%  & 0.4342 & $0.64\times$ & 16.0\% \\
8 & 0.2056 & $1.67\times$ & 20.9\%  & 0.4881 & $0.46\times$ & 5.8\% \\
\bottomrule
\end{tabular}
\label{tab:scaling_table}
\end{table}

\begin{table}[H]
\centering
\caption{Thread Scaling for Depth-Limited Quick Sort ($N = 10,000$, Initial Depth = 8)}
\vspace{4pt}
\begin{tabular}{@{}ccccc@{}}
\toprule
\textbf{Threads ($P$)} & \textbf{Sequential Time (s)} & \textbf{Parallel Time (s)} & \textbf{Speedup ($S$)} & \textbf{Parallel Efficiency ($E$)} \\
\midrule
1 & 0.001783 & 0.003939 & $0.45\times$ & 45.0\% \\
2 & 0.001728 & 0.003960 & $0.44\times$ & 22.0\% \\
4 & 0.002134 & 0.006568 & $0.32\times$ & 8.0\% \\
8 & 0.002406 & 0.002871 & $0.84\times$ & 10.5\% \\
\bottomrule
\end{tabular}
\label{tab:depth_table}
\end{table}

\section{Detailed Comparative Analysis and Discussion}

\subsection{OpenMP Sections vs. OpenMP Tasks}
OpenMP provides multiple paradigms to exploit irregular and recursive concurrency:
\begin{itemize}
    \item \textbf{Parallel Sections (\texttt{\#pragma omp sections})}: Represents static task parallelism. Each section denotes a static block assigned to an available thread in the current team. When combined with recursive depth limits (e.g., in \texttt{parallel\_merger\_sort.c}), the number of concurrently active sections is strictly bounded by $2^{\text{depth}+1}$. This caps thread contention and yields positive speedup ($1.65\times$ on 4 threads).
    \item \textbf{Explicit Tasks (\texttt{\#pragma omp task})}: Designed for dynamic and irregular parallelism. Threads dynamically extract tasks from a centralized or work-stealing task pool. While conceptually ideal for recursive divide-and-conquer, spawning explicit tasks incurs non-negligible cost: task context allocation on the heap, environment variable capture, queue synchronization, and thread scheduling.
\end{itemize}

\subsection{Granularity and the Task Creation Overhead Catastrophe}
The critical distinction between \texttt{parallel\_merger\_sort.c} and \texttt{parallel\_quick_sort.c} lies in \textbf{granularity control}. In \texttt{parallel\_quick\_sort.c}, tasks are spawned recursively until $low \ge high$:
\[
\text{Total Tasks Generated} = 2n - 1 \approx 2 \times 10^6 \text{ tasks for } N = 10^6
\]
For minuscule sub-arrays (e.g., sub-arrays of size $2$ or $3$), the time required to sort the array sequentially is mere nanoseconds (a handful of CPU cycles), whereas allocating, queueing, and context-switching an OpenMP task requires hundreds of nanoseconds. Consequently, as the thread count increases from 1 to 8 (Table~\ref{tab:scaling_table}), thread contention on the OpenMP runtime task queues escalates, inflating parallel execution time from 0.310s to 0.488s and causing severe performance degradation (speedup drops from $0.83\times$ to $0.46\times$).

\subsection{Impact of Depth Limits and Base-Case Cutoffs}
In \texttt{quicksort_depth.c}, task explosion is prevented by establishing a recursion depth threshold $depth = 8$. This limits the maximum number of generated tasks to:
\[
\sum_{k=0}^{7} 2^k = 2^8 - 1 = 255 \text{ tasks}
\]
which adequately saturates modern multi-core processors without overloading the runtime task scheduler. 

Furthermore, upon hitting the depth cutoff, switching to \texttt{insertionSort()} eliminates recursive call overhead. For small contiguous partitions, Insertion Sort performs remarkably well because it executes entirely in the L1 CPU data cache and requires zero dynamic memory allocations. However, for a small total workload ($N=10,000$ elements, total execution time $\approx 1.3\text{ ms}$), the fixed overhead of thread team creation and barrier synchronization still outweighs parallel gains ($S = 0.31\times$ to $0.84\times$), demonstrating that parallel sorting requires sufficiently large problem sizes ($N \ge 10^5$) to amortize OpenMP runtime startup costs.

\subsection{Memory Bandwidth and In-Place Sorting Trade-offs}
An additional critical architectural factor observed is memory traffic:
\begin{itemize}
    \item \textbf{Merge Sort}: Requires dynamic allocation of an auxiliary buffer (\texttt{int *temp = malloc(...)}) during every merge step. While the merge step is sequential at the bottom levels, repeated memory allocation and data copying from \texttt{temp} back to \texttt{a} places substantial demand on the memory bus, bounding parallel efficiency to $\approx 41\%$.
    \item \textbf{Quick Sort}: Operates completely in-place via element swapping. Once partitioned, sub-arrays reside in localized cache lines. If partitioned with adequate granularity, in-place Quick Sort avoids memory allocation bottlenecks.
\end{itemize}

\section{Conclusion}
The empirical results and theoretical analysis demonstrate the following core findings:
\begin{enumerate}
    \item Parallel divide-and-conquer cannot be naively implemented by wrapping every recursive branch in an OpenMP task or section. Doing so results in parallel slowdown due to task scheduling and queue management overheads exceeding the payload computation.
    \item Restricting recursion depth (as in \texttt{parallel\_merger\_sort.c} with $\text{depth} \le 4$ and \texttt{quicksort_depth.c} with $\text{depth} \le 8$) effectively limits task proliferation and ensures that parallel threads execute sufficiently coarse chunks of work.
    \item Switching to an efficient cache-friendly algorithm such as Insertion Sort for base-case sub-arrays is vital for optimizing memory locality and minimizing recursion stack overhead.
    \item Parallel sorting yields meaningful speedups only when the input dataset is sufficiently large to offset thread creation, synchronization, and synchronization barrier overheads.
\end{enumerate}

\end{document}
